\section{Reverse excision and homology localizers}
\label{sec:A}
We first compare the actual span and sieve axioms, then test the proposed implication to $RT$-fixedness. The derivator assertions require a separate analysis: an axiom on squares of coefficient objects is stronger than an axiom on the cohomological localizer.

\subsection{The precise comparison of the excision axioms}

\begin{claim}{A.1}{Proved}{Specified span comparisons}\label{A1}
Let $W$ be a fundamental localizer, with the relative Theorem A. Put $\Psi=(a\leftarrow0\to b)$ and $\Phi=\Psi^{\op}$. For a strict functor $F:\Psi\to\Cat$, write $T_cF$ for its covariant Grothendieck construction and $T_fF$ for the fibred construction of $F:\Phi^{\op}\to\Cat$. The canonical maps from the outer fibers to these two total categories represent the two outer maps of the same homotopy pushout of spaces. All these comparisons are natural in a strict morphism $F\to F'$, including the maps from the middle fiber up to the specified natural transformations. For a cover $X=U\cup V$ by full sieves, the augmentation $q:T_c(U\leftarrow U\cap V\to V)\to X$ has a canonical fully faithful left adjoint $\alpha:X\to T_c(U\leftarrow U\cap V\to V)$, and both are weak homotopy equivalences.
\end{claim}
\begin{proof}
An object of $T_cF$ is $(j,x)$ with $x\in F(j)$. Its arrows are pairs $(\alpha,h):(j,x)\to(k,y)$ with $\alpha:j\to k$ in $\Psi$ and $h:F(\alpha)x\to y$ in $F(k)$. Composition is $(\beta,l)(\alpha,h)=(\beta\alpha,lF(\beta)h)$. The full subcategories $M_a=F(0)\cup F(a)$ and $M_b=F(0)\cup F(b)$ are sieves, cover $T_cF$, and intersect in $F(0)$. The functor $r_a:M_a\to F(a)$ acts as $F(0\to a)$ and the identity on the two layers. It is left adjoint to the inclusion $s_a:F(a)\to M_a$, with identity counit and crossing-arrow unit. The same formulas define $r_b,s_b$. Thus
\[
 (M_a\leftarrow F(0)\to M_b)
 \xrightarrow{(r_a,\id,r_b)}
 (F(a)\leftarrow F(0)\to F(b))
\]
is an actual map of spans which is objectwise a nerve homotopy equivalence. The nerve of the cover is the union of the two nerves: a composable chain lies in the sieve containing its last object. The inclusions of the intersection are simplicial monomorphisms. Consequently its realization is a homotopy pushout, and the displayed map induces the required weak equivalence with the double mapping cylinder. The natural transformations $s_ar_a\Leftarrow\id_{M_a}$ and $s_br_b\Leftarrow\id_{M_b}$ give the two homotopies from the middle-fiber inclusion, not a strictly commutative square in $\Cat$.

For the fibred convention, $T_fF=(T_c(F^{\op}))^{\op}$, where $F^{\op}:\Psi\to\Cat$ has fibers $F(j)^{\op}$. There is a natural homeomorphism $B(C^{\op})\to BC$: reverse the order of each chain and its barycentric coordinates. Apply this both to the total category and to its fibers. The preceding double mapping cylinder comparison then identifies the prescribed outer-fiber maps for $T_fF$ with the same homotopy pushout. If actual categorical zigzags are desired, apply singular simplices and their simplex-category functor to this natural diagram of spaces. The counits $|\operatorname{Sing}Z|\to Z$ and the subdivision comparisons $B(\Delta/K)\to|K|$ are natural weak equivalences. This specifies the arrows and their directions; the comparison is not an asserted equality of categories. The general Grothendieck comparison is also \cite[Theorem 26.8 and Corollary 26.9]{AChacholskiScherer}.

In the cover case $q$ forgets the layer tag and is induced by the three inclusions into $X$. The natural double mapping cylinder map onto $BU\cup_{B(U\cap V)}BV=BX$ is a weak equivalence, since both attachments are CW subcomplex inclusions. It factors through the comparison just given and $Bq$. This proves the assertion, including empty intersections and empty pieces. By minimality and saturation of fundamental localizers, every comparison lies in $W$; the exact general statements used are \cite[Proposition 4.2.4 and Corollary 4.2.19]{Cisinski2006B}.
\end{proof}

We can also recover the actual arrow $\alpha:X\to\overline X$ mentioned in cote 105/045. Define $\phi:X\to\Psi$ by $\phi(x)=0$ for $x\in U\cap V$, $\phi(x)=a$ for $x\in U\setminus V$, and $\phi(x)=b$ for $x\in V\setminus U$. This is a functor: if $x\to y$ and $y$ lies in either sieve, so does $x$, and therefore the unique permitted arrow $\phi(x)\to\phi(y)$ exists. Put
\[
 \alpha(x)=(\phi(x),x),\qquad
 \alpha(h:x\to y)=(\phi(h),h).
\]
For $(j,y)$ in the total category, every arrow $x\to y$ in $X$ forces $\phi(x)\to j$, and that base arrow is unique. Consequently
\[
 \Hom_{T_cF}(\alpha x,(j,y))\cong\Hom_X(x,q(j,y))
\]
naturally. Hence $\alpha\dashv q$, with $q\alpha=\id_X$ and counit $\alpha q(j,y)\to(j,y)$ the crossing arrow on $\id_y$. This is a direct proof that both maps lie in $W$. It includes empty pieces and an empty intersection. The construction of $\alpha$ is strictly natural for maps over the classifying functor $\phi:X\to\Psi$. A general cover-preserving map can send an exclusive stratum into the intersection; then it has the canonical comparison $\alpha' h\Rightarrow (T_ch)\alpha$, while $q$ remains strictly natural. Thus the manuscript's direction is valid, and is distinct from its right-adjoint augmentation. Our first checkpoint's assertion that no canonical intersection tag was available was too strong; the tag $0$ is canonical and gives the displayed adjunction.

\begin{claim}{A.2}{Proved}{The cofiber criterion, without accessibility}\label{A2}
For a fundamental localizer $W$, reverse sieve excision, XII W(7 bis), and
\[
 f\in W\quad\Longleftrightarrow\quad C(Bf)\text{ is }W\text{-aspherical}
 \tag{A.2.1}
\]
are equivalent. The forward sieve implication is automatic. Statements about a space here mean statements about its singular simplex category, compared naturally with its nerve. Under these conditions the pointed aspherical class is closed under pointed homotopy colimits, retracts, and all three implications of two-out-of-three in pointed cofiber sequences.
\end{claim}
\begin{proof}
We give the categorical cofiber model to avoid replacing a map by an abstract homotopy type. For $f:A\to B$, form the collage $M_f$ with layers $A,B$, original internal arrows, crossing arrows
\[
 \Hom_{M_f}(a,b)=\Hom_B(fa,b),\qquad
 \Hom_{M_f}(b,a)=\varnothing.
\]
Adjoin a point $t$ with a unique arrow $a\to t$ for each $a\in A$ and no arrows between $B$ and $t$. Call the result $K(f)$. The subcategories $M_f$ and $A^{\triangleright}$ are full sieves covering $K(f)$ and intersecting in $A$. The retraction $r:M_f\to B$, with $r|_A=f$ and $r|_B=\id$, is left adjoint to $B\hookrightarrow M_f$. Claim A.1 identifies $BK(f)$ with $C(Bf)$, at the specified cone point. A commutative square $(u,v):f\to f'$ gives the functor $K(f)\to K(f')$ by $u,v$ on the layers and $t\mapsto t'$.

We also record the forward gluing proof. For $\eta:F\to G$ over any small category $J$, with every $\eta_j\in W$, the evaluation functor $(T_cF)/j\to F(j)$, $(i,x,\alpha:i\to j)\mapsto F(\alpha)x$, is left adjoint to $y\mapsto(j,y,\id)$. Its unit and counit give nerve equivalences natural in $F$. Relative Theorem A therefore gives $T_c\eta\in W$. Applying this to spans and then A.1 proves forward gluing and its compatibility with cofiber models. This is the step rechecked from DER\_FOLLOWUP A.10; it uses neither accessibility nor an interchange of $\Ho(\Cat^J)$ with $\Ho(\Cat)^J$.

If $f\in W$, compare the span $(B\leftarrow A\to e)$ with $(B\leftarrow B\to e)$. Forward gluing gives $K(f)\to K(\id_B)\in W$, and the latter category has contractible nerve. Conversely, if $K(f)$ is aspherical and reverse sieve excision holds, then $A^{\triangleright}\to K(f)$ is in $W$. In its sieve-cover square this implies $A\to M_f\in W$, whence $f=r\circ(A\to M_f)\in W$.

Assume (A.2.1). For a homotopy pushout, the cofibers of parallel arrows are naturally weakly equivalent, by associativity of the pushout after adjoining the cone point. Taking the iterated double mapping cylinder, or the iterated Grothendieck construction followed by its augmentation, gives the actual comparison. Thus (A.2.1) proves both sieve implications. The empty source is included: $K(\varnothing\to B)=B\amalg e$; for the empty identity it is $e$.

For pointed spaces, a homotopy colimit of aspherical objects is aspherical by the Grothendieck argument applied to the pointed diagram and the diagram of points. This uses pointed homotopy colimits: the section of cone points in an unpointed replacement is collapsed. Retracts follow from the retract closure of $W$. The cofiber sequence $X\to *\to\Sigma X$ and (A.2.1) show that $X$ is aspherical if and only if $\Sigma X$ is. If two vertices in $X\to Y\to Z$ are aspherical, forward gluing proves the missing vertex when it is $Z$; cofiber rotation and desuspension give the other two implications. These statements use the basepoint section $*\to X$ and two-out-of-three to replace aspherical spaces by points. Every nonempty pointed realization has the required section. No basepoint is chosen in an empty unpointed source.
\end{proof}

\begin{claim}{A.3}{Proved}{The strengthened cube adds no condition}\label{A3}
For a fundamental localizer the following are equivalent:
\begin{enumerate}[label=(\roman*),leftmargin=2.2em]
\item cote 105 W7 bis for every strict span;
\item XII W(7 bis), equivalently XVI Loc (6$'$ a,b);
\item cote 105 W8 bis for every strict morphism of spans;
\item XII W(7), equivalently XVI Loc (6 a,b,c).
\end{enumerate}
These equivalences hold for arbitrary, as well as accessible, fundamental localizers. The coefficient in W8 bis condition (a) is the total functor $f$, not $f_a$.
\end{claim}
\begin{proof}
Claim A.1 transfers a sieve square to its span total category and transfers an arbitrary span to the sieve square $(M_a,M_b)$ inside its total category. The retractions $r_a,r_b$ and the total comparison are in $W$. Two-out-of-three therefore transfers each edge test in exactly the required direction. This proves (i)$\Leftrightarrow$(ii), with the fibred convention handled by the same natural zigzag. Applying the construction to a strict morphism of spans transfers all four vertical maps, proving (iii)$\Leftrightarrow$(iv).

Here is the implication from square to cube. For a morphism $f_i:X_i\to Y_i$ of two homotopy pushouts, let $C_i=C(Bf_i)$, $i=0,a,b,1$. Adding a disjoint basepoint and then taking cofibers are colimit functors, so the $C_i$ form a homotopy pushout in pointed spaces. The comparison is the colimit-interchange map on the expanded diagram, not an abstract isomorphism. If $C_0,C_a,C_b$ are aspherical, so is $C_1$ by forward gluing. If $C_a,C_b,C_1$ are aspherical, the map of spans to $(*\leftarrow C_0\to*)$ gives $C_1\to\Sigma C_0$ in $W$. Hence $\Sigma C_0$, and then $C_0$, are aspherical. Finally, if $C_0,C_1$ are aspherical, the basepoint map $*\to C_0$ is a $W$-map. The induced map $C_a\vee C_b\to C_1$ is in $W$. Each summand is a pointed retract of the wedge, so each is aspherical. By A.2 these are exactly the three cube implications.

Conversely, the cube axiom implies the reverse square axiom by comparing the span $(X_a\leftarrow X_0\to X_b)$ with $(X_a\leftarrow X_a\to P)$, where $P$ is its homotopy pushout, using the given maps and the coherent square. Equivalently, for a sieve square one uses the degenerate-cover cube in XII Proposition 6: the three vertical tests reduce to the opposite-edge tests. To avoid any strictification ambiguity, the same conclusion follows from the $C_i$ construction applied to the map from a span to its constant total span and then the collage replacement in A.2; all added comparisons are weak equivalences. Thus full cube excision entails A.2. One may also read the elementary degenerate cube explicitly as the morphism from $(U\cap V\to U,\,U\cap V\to V)$ to $(V\to X,\,V\xrightarrow{\id}V)$: its maps are $U\cap V\to V$, $U\to X$, and $\id_V$, while its total map is a weak equivalence. The reverse cube implication gives the required missing edge.

The text reads ``Alors deux de ces conditions entraînent la troisième.'' If (a) were $f_a\in W$, (a) would be part of (c), and (a)+(c) would force $f_0\in W$ regardless of the total map. This is false even for integral homology: take the strict span morphism $(e\leftarrow\{0,1\}\to e)\to(e\leftarrow e\to e)$ of discrete finite categories. Both outer maps are identities while the middle map sends $\Z^2\to\Z$ by addition on $H_0$ and is not a homology equivalence. The program \texttt{code/A\_span\_glyph.py} verifies these functors and their four- and three-object total categories. Thus the mathematics forces the reading $f\in W$ if the intended assertion is a valid homological cube axiom. It cannot certify the uncertain ink. The uncertain margin concerning a cartesian map of total categories is consistent with the equivalence just proved for maps preserving the specified transport; it supplies no stronger axiom for arbitrary maps over a base.
\end{proof}

\begin{claim}{A.4}{Proved}{What reversibility does force}\label{A4}
Every reversible fundamental localizer contains $W_{\Z}$. The class $W_{\Z}$ is itself reversible and is the smallest such localizer. This does not imply that every reversible localizer is $RT$-fixed.
\end{claim}
\begin{proof}
For an integral homology equivalence $f$, the actual map on free $H_0$ groups is induced by the component function, so its invertibility forces that function to be bijective, even for infinite sets. If its source is empty, both categories are empty. Otherwise $C(Bf)$ is connected and integrally acyclic. Its suspension is simply connected and acyclic; Hurewicz applied to a first possible nonzero homotopy group, followed by Whitehead, makes it weakly contractible. Two applications of A.2 therefore put $f$ in $W$. Integral homology has the cofiber criterion by its reduced long exact sequence, which proves reversibility of $W_{\Z}$. Its relative Theorem A follows by applying chains to the natural comma-category homotopy-colimit decomposition in A.2. These arguments recheck POLES B.5--B.6 at the precise points used.
\end{proof}

\subsection{The sphere-null counterexample and the stable witness}

\begin{claim}{A.5}{Proved}{A complete membership test and all localizer axioms}\label{A5}
Let
\[
 \Nul_2=\{X\in\cS_*:\Map_*(X,S^n_{(2)})\simeq *\text{ for every }n\geq2\},
 \qquad W_S=\{f:A\to B:C(Bf)\in\Nul_2\}.
\]
Mapping spaces include all components. Then $W_S$ is a reversible fundamental localizer and satisfies all four equivalent axioms in A.3. It contains $B C_2\to e$ and does not contain $P\to e$, where $P$ is the seven-object face-occurrence category of the usual two-cell presentation of $\RP^2$.
\end{claim}
\begin{proof}
The classical $2$-local EHP fiber sequence is
\[
 S^n_{(2)}\longrightarrow\Omega S^{n+1}_{(2)}
              \longrightarrow\Omega S^{2n+1}_{(2)},\qquad n\geq1;
\]
we use \cite[Theorem 2.2, with its preceding localization convention]{ABoyde} and \cite[slide 2 and Theorem 1(i),(ii)]{ARavenelEHP}. If $\Sigma X\in\Nul_2$, mapping $X$ into this sequence makes both loop-space terms contractible for every $n\geq2$, so the fiber is contractible. This proves desuspension closure, including $n=2$ and components. Mapping out of pointed homotopy colimits gives homotopy limits, so $\Nul_2$ is closed under such colimits, weak equivalences, and retracts. Cofiber rotation, the mapping-space fiber sequence, and desuspension closure give all three cofiber implications. It contains the point. Its members are connected: if a second component existed, evaluation on it would exhibit $S^2_{(2)}$ as a retract of $\Map_*(X,S^2_{(2)})$.

For completeness we verify the localizer axioms directly. The cofiber of an identity is contractible. For composable unpointed maps $f,g$, cofiber pasting gives a pointed cofiber sequence $C(f)\to C(gf)\to C(g)$; this can be constructed as the successive quotients of the two mapping cylinders. Its two-out-of-three implies that for $W_S$. Cofiber takes an arrow retract to a pointed homotopy retract, proving retract closure. To verify Loc (1)c without assuming saturation in advance, we identify $W_S$ as a mapping-equivalence class. If $C(f)$ is null, the rotated sequence makes $\Map(f,\Omega S^j_{(2)})$ an equivalence for every $j\geq2$: exactness against the next cofiber term gives surjectivity on components, and translation in loop spaces identifies all its homotopy fibers with the contractible fiber at the constant map. Comparing the EHP fiber sequences makes $\Map(f,S^n_{(2)})$ an equivalence. The converse follows by taking the fiber at the constant map. Thus $W_S$ is exactly the class inverted by all $\Map(-,S^n_{(2)})$, and identities, two-out-of-three, Loc (1)c and strong saturation follow immediately. A.7 supplies a separate telescope proof of Loc (1)c for the more general cofiber constructions used later.

A category with a final object has contractible nerve, proving Loc (2). For $A\xrightarrow{u}B\to J$, the natural weak equivalences
\[
 \hocolim_{j\in J}B(A/j)\longrightarrow BA,\qquad
 \hocolim_{j\in J}B(B/j)\longrightarrow BB
\]
are the Grothendieck augmentations: the projection to $A$ has the section $a\mapsto(a,\id_{v(a)})$ as an adjoint. Their naturality in $u$ is literal. If every $u/j$ is in $W_S$, the pointed homotopy colimit of its cofibers is $C(Bu)$, since $(-)_+$ and cofiber preserve colimits. Hence $u\in W_S$, proving relative Theorem A. Finally a nonempty pointed $Y$ satisfies $Y\to *\in W_S$ exactly when $\Sigma Y\in\Nul_2$, equivalently $Y\in\Nul_2$. Thus A.2 applies and A.3 verifies all the requested reverse axioms. For $\varnothing\to B$ with $B\ne\varnothing$, the cofiber $BB_+$ is disconnected and fails the test; the empty identity passes.

Miller's \cite[Theorem C, p.~40]{AMiller} applies to each nilpotent $S^n_{(2)}$, whose mod-$2$ homology is bounded. It gives $\RP^\infty=B C_2\in\Nul_2$. On the other hand obstruction theory for a two-dimensional source and a simply connected target gives
\[
 [\RP^2,S^2_{(2)}]_*\cong H^2(\RP^2;\Z_{(2)})\cong\Z/2.
\]
The last equality follows from the cellular cochain map $\Z_{(2)}\xrightarrow{2}\Z_{(2)}$. The nonzero class excludes $\RP^2$. The finite categorical model and its integer chain matrices are constructed and checked by \texttt{code/A\_finite\_categories.py}; it has seven objects, 25 arrows, normalized chain ranks $(7,18,12)$ and the single nonunit Smith factor $2$. Its geometric realization is the face subdivision of the disk attached to a circle by degree two. This proves the two asserted membership tests for actual functors in $\Cat$.
\end{proof}

\begin{claim}{A.6}{Proved}{The named witness in the stable closure}\label{A6}
With $V=\Sigma^\infty\RP^\infty$ and $\mathbb S/2$ the mod-$2$ Moore spectrum,
\[
 \loc(V)=\loc(\mathbb S/2).
\]
Consequently $P\to e\in R(T(W_S))\setminus W_S$. In particular the assertion in Part A.1 is false.
\end{claim}
\begin{proof}
Every localizing subcategory of spectra is a tensor ideal: for $E$ in it, the spectra $F$ with $E\wedge F$ in it form a localizing subcategory containing the sphere, and sphere cells generate all spectra. Let $L=\loc(V)$. The Snaith splitting \cite[Theorem 1.1 and the filtration in Section 2]{ASnaith} identifies
\[
 \Sigma^\infty Q\RP^\infty\simeq
 \bigvee_{j\geq1}V^{\wedge j}_{h\Sigma_j}.
\]
All terms are in $L$ by its tensor and homotopy-colimit closure. The space-level Kahn--Priddy section, \cite[unnumbered main theorem, p.~103, and its construction on p.~110]{AKahnPriddy}, makes $(Q_0S^0)_{(2)}$ a homotopy retract of $(Q\RP^\infty)_{(2)}$. This is a section of spaces; an infinite-loop section is not asserted. Suspension spectra of nilpotent space localizations agree with localizations of their suspension spectra, because the canonical comparison is an integral homology equivalence between connective spectra. Thus $\Sigma^\infty(Q_0S^0)_{(2)}\in L$.

For a connected grouplike infinite-loop space $Z$, the bar construction $BZ$ is the realization of the simplicial space of finite powers of $Z$. Stably the reduced suspension spectrum of a finite product splits into the finite wedge of nonempty smash products. Hence $\Sigma^\infty Z\in L$ implies $\Sigma^\infty BZ\in L$. Apply this successively to $Z=\Omega^\infty E$, $E=\tau_{\geq1}\mathbb S_{(2)}$. All its iterated deloopings remain connected, and
\[
 E\simeq\hocolim_{r\geq0}\Sigma^{-r}\Sigma^\infty\Omega^\infty\Sigma^rE.
\]
The comparison is induced by the counits of stabilization. For each fixed homotopy degree its connectivity tends to infinity, by Freudenthal applied to the increasingly connected infinite-loop spaces. Thus $E\in L$. The identical bar argument starting from $B C_2$ gives $H\F_2\in L$. Smashing $E\to\mathbb S_{(2)}\to H\Z_{(2)}$ with $\mathbb S/2$ now gives a triangle whose first and third terms belong to $L$ and whose middle term is $\mathbb S/2$.

Conversely $V[1/2]$ is connective and has zero integral homology, and is therefore zero by stable Hurewicz. The kernel of $(-)[1/2]$ is $\loc(\mathbb S/2)$: the cofibers $\mathbb S/2^r$ are finite extensions of $\mathbb S/2$, and the fiber of $F\to F[1/2]$ is the telescope of $\Sigma^{-1}F/2^r$. Tensor closure proves the claim for every spectrum $F$. Hence $V\in\loc(\mathbb S/2)$, proving equality.

The cofiber of $B C_2\to *$ is $\Sigma\RP^\infty$, so $T(W_S)$ contains $L$. It therefore contains $\Sigma^\infty C(BP\to*)\simeq\Sigma^2\mathbb S/2$. The definition of $R$ places $P\to e$ in $R(T(W_S))$, while A.5 excludes it from $W_S$. This checks the Galois argument in POLES F.2--F.6 by its actual generators: $T(W)\subseteq L$ if and only if every $\Sigma^\infty C(Bf)$ for $f\in W$ belongs to $L$, if and only if $W\subseteq R(L)$.
\end{proof}

\begin{claim}{A.7}{Proved}{An accessible counterexample in ZFC}\label{A7}
There exists an accessible fundamental localizer satisfying W7 bis, W8 bis, XII W(7 bis), XII W(7), and XVI Loc (6 a,b,c), which is not $RT$-fixed. One fully specified example has membership
\[
 \begin{split}
 f\in W_c\quad\Longleftrightarrow\quad&
 \text{every countable pointed subcomplex of a CW model}\\
 &\text{of }C(Bf)\text{ lies in a countable }\Nul_2\text{-subcomplex}.
 \end{split}
\]
The condition is independent of the CW model. In particular accessibility does not rescue Part A.1.
\end{claim}
\begin{proof}
We recheck the bounded-subcomplex argument of REVLOC B.1--B.4. For any pointed class $\mathcal A$ closed under weak equivalences, pointed homotopy colimits, and cofiber two-out-of-three, define $\mathcal A_c$ by the countable enlargement condition in the statement. If $X\in\mathcal A_c$, its countable $\mathcal A$-subcomplexes form a filtered poset with union $X$. Their inclusions are cofibrations and every compact sphere or homotopy meets a finite subcomplex, so this union is also their homotopy colimit. Thus $\mathcal A_c\subseteq\mathcal A$.

A cellular map sends a countable subcomplex into a countable subcomplex, since a compact characteristic disk meets finitely many cells; the same is true of the tracks of a cellular homotopy. Given a cellular homotopy equivalence $X\rightleftarrows Y$, start with a countable subcomplex of $Y$, close alternately under both maps and both inverse homotopies, and enlarge the $X$ side to a countable $\mathcal A$-subcomplex at each stage. The two countable unions are homotopy equivalent; the $X$ union is a sequential homotopy colimit of members of $\mathcal A$. The $Y$ union therefore belongs to $\mathcal A$. This proves model independence.

For a pointed homotopy colimit use the cellular bar model. Any countable collection of its cells uses countably many objects and arrows of the indexing category, which generate a countable subcategory. Enlarge the specified cells in each involved vertex to a countable $\mathcal A$-subcomplex, close under all these structure maps, and repeat countably. The resulting vertexwise subdiagram is countable and in $\mathcal A$ at each vertex. Its pointed bar realization is a countable $\mathcal A$-subcomplex containing the original cells. This proves homotopy-colimit closure of $\mathcal A_c$.

For cofibers replace the map by a cellular cofibration $X\subset Y$, with quotient $C=Y/X$. A countable subcomplex $Y_0$ determines $X_0=Y_0\cap X$ and $C_0=Y_0/X_0$. A prescribed countable subcomplex of $C$ can be lifted by taking its nonbasepoint cells in $Y\setminus X$ and their boundaries. If any two of $X,Y,C$ lie in $\mathcal A_c$, enlarge those two vertices to countable $\mathcal A$-subcomplexes, restore these compatibility equations, and iterate. The unions form a cofiber sequence; two vertices belong to $\mathcal A$ by sequential closure, and the third does by cofiber two-out-of-three. This proves the same property for $\mathcal A_c$. Retracts follow by homotopy-idempotent telescopes. For clarity, the general cofiber construction also satisfies the precise, two-composite version of Loc (1)c), independently of the special mapping-space test in A.5. Indeed it is closed under homotopy colimits of arrows, because cofibers commute with those colimits. If $f:X\to Y$ and $g:Y\to X$ have $gf,fg$ in the class, form $X_\infty=\operatorname{Tel}(gf)$ and $Y_\infty=\operatorname{Tel}(fg)$. The inclusion $X\to X_\infty$ is in the class: it is the homotopy colimit of the objectwise equivalences $(gf)^n$ from the constant $X$ diagram to this telescope diagram. The same holds for $Y$. The composites of the induced maps $f_\infty,g_\infty$ are the self-maps induced by the telescope structure maps. Those self-maps are homotopy equivalences: shifting the stage index gives their homotopy inverses. Therefore $f_\infty,g_\infty$ are weak equivalences. The commuting square with the two inclusions gives $f,g$ in the class by two out of three. Identities, arrow retracts, terminal objects, and relative Theorem A are exactly the cofiber and Grothendieck-construction arguments of A.5, which used only the closure properties already verified here. The cofiber criterion gives reversibility.

There is a set of countable CW complexes and maps up to homotopy-model choices (their underlying sets have cardinal at most the continuum). Take all arrows between countable such models whose cofibers belong to $\mathcal A$. Their simplex-category representatives form a small set $S$. If a map has cofiber in $\mathcal A_c$, replace it by its mapping-cylinder inclusion $X\subset M$. Every countable subpair lies in a countable subpair $X'\subset M'$ whose quotient is a countable $\mathcal A$-subcomplex of $M/X$: lift the noncollapsed cells of the chosen enlargement and their attaching boundaries. The quotient is exactly that subcomplex. These subpairs form a filtered system with homotopy colimit the original arrow. Relative Theorem A, via its Grothendieck construction, puts the arrow in every fundamental localizer containing $S$. Conversely each generator belongs to the displayed cofiber localizer. Thus it is the fundamental localizer generated by $S$, proving accessibility. If a nonempty $Y_+$ were in $\mathcal A_c$, then its retract $S^0$ would be there; all pointed spaces would follow by cells, and the resulting total localizer is also accessible, generated by $\varnothing\to e$. Thus this edge case creates no omitted exception.

Now take $\mathcal A=\Nul_2$. The countable complex $\RP^\infty$ lies in $\mathcal A_c$, while $\RP^2$ does not even lie in $\mathcal A$. Claim A.6 supplies the same stable witness for $W_c$. This proves the assertion without requiring the stronger claim that the entire $W_S$ is accessible.
\end{proof}

\begin{claim}{A.8}{Refuted}{The manuscript's assertion for $W_\infty$}\label{A8}
XII Théorème 2, ms.~78, is false, and the homological-only inference appended to R3 is also false. The first fails for the explicit 37-object category below; the second fails by A.5--A.7. The corrected statement is: the four localizer axioms in A.3 all express reversibility, force $W_\Z\subseteq W$, and do not force $RT$-fixedness, even under accessibility.
\end{claim}
\begin{proof}
Take the presentation complex
\[
 L=\langle a,b\mid a^5b^{-3},\ b^3(ab)^{-2}\rangle.
\]
Its cellular boundary $d_2=\left(\begin{smallmatrix}5&-2\\-3&1\end{smallmatrix}\right)$ has determinant $-1$, so $L$ is connected and integrally acyclic. The permutations $a=(1\ 4\ 3\ 2\ 5)$ and $b=(1\ 2\ 3)$ satisfy both relators and $a\ne1$, proving $\pi_1L\ne1$. Cone-triangulate the two relator disks and take the face-occurrence category $A$: an arrow to a simplex is an increasing injection selecting that occurrence of its source face, and composition is composition of injections. This has 35 objects and 159 arrows; its nerve realizes the subdivision of $L$.

Adjoin two cone points $t,s$, each receiving one arrow from every object of $A$, and no arrows between them. In $X=A\star\{t,s\}$ the full sieves $U=A\star\{t\}$ and $V=A\star\{s\}$ cover $X$. Both pieces are contractible; $BX$ is the unreduced suspension of the connected acyclic $BA$, hence is simply connected and acyclic and therefore weakly contractible. The intersection $A$ is not. Thus $U\to X$ belongs to $W_\infty$ but $A\to V$ does not. The map of the cover square to the constant point square also refutes the reverse cube implication.

The included program reconstructs all arrows and composites, verifies all 1,813 associativity triples for $X$, both sieve conditions, the relators in $A_5$, and the integral nerve-chain complexes. The normalized chain ranks for $A$ are $(35,124,90)$ and for $X$ are $(37,194,338,180)$; every nonzero Smith factor is $1$, giving the stated homology. These computations recheck TOTASPH A.13 and POLES B.7--B.8. The proof of noncontractibility uses the explicit nontrivial quotient of the presentation group; it is not inferred from a numerical computation of homology.
\end{proof}

\begin{claim}{A.9}{Proved}{Audit of the full torsion-finite realization in REVLOC D.2.2}
\label{A9}
Let $\mathcal C$ be any nonzero thick subcategory of torsion finite spectra and put $L=\loc(\mathcal C)$. There is an accessible fundamental localizer $W$ satisfying every equivalent reverse axiom of A.3, with $T(W)=L$ and $W\ne RT(W)$. The following construction specifies its membership for all small-category functors and a separating functor. Thus the conclusion of REVLOC D.2.2 survives the strengthened span axioms, at every prime.
\end{claim}
\begin{proof}
Choose a prime $p$ and a nonzero $p$-primary $F\in\mathcal C$. Primary summands exist and remain in $\mathcal C$: the endomorphism group of a torsion finite spectrum is finite, and the Chinese-remainder idempotents split in the finite stable category. For $p=2$ use $\Nul_2$. For odd $p$, replace the sphere targets by the collection
\[
 \mathcal Z_p=\{S^{2n-1}_{(p)},\ J_{p-1}(S^{2n})_{(p)}:n\geq1\},
 \qquad
 \Nul_p=\{X:\Map_*(X,Z)\simeq *\ (Z\in\mathcal Z_p)\}.
\]
Here $J_{p-1}$ is the indicated finite James stage. The two $p$-local Toda fiber sequences are
\begin{align*}
 J_{p-1}(S^{2n})_{(p)}&\longrightarrow\Omega S^{2n+1}_{(p)}
                        \longrightarrow\Omega S^{2np+1}_{(p)},\\
 S^{2n-1}_{(p)}&\longrightarrow\Omega J_{p-1}(S^{2n})_{(p)}
                        \longrightarrow\Omega S^{2np-1}_{(p)}.
\end{align*}
We inspected the printed formulas in \cite[p.~81, proof (ii)]{ABodigheimerHenn}. They imply desuspension closure by exactly the mapping-fiber argument of A.5. The bottom case includes the circle and the finite James stage on $S^2$; no even-sphere fiber is substituted at an odd prime. These targets are nilpotent and have bounded mod-$p$ cohomology, so Miller's Theorem~C makes $BC_p$ null for every target. The closure and localizer proofs of A.5 and A.7, including the telescope verification of Loc~(1)c), therefore apply to $\Nul_p$.

We check the infinite-loop input in the earlier proof against a stronger opened source. \cite[Theorem~1.5]{ALevy} identifies the classes generated under colimits and extensions by $\Sigma^\infty BC_p$ and $\Sigma\mathbb S/p$. Its Proposition~3.1 applies Bousfield's \cite[Proposition~3.1]{ABousfield1996} to pass from connective spectral nullification to ordinary nullification of infinite-loop spaces. Consequently \cite[Theorem~3.2]{ALevy} says that every $\Omega^\infty E$, with $E$ connected and $E[1/p]=0$, is ordinarily $BC_p$-acyclic. Since the reversible cofiber localizer of $\Nul_p$ contains $BC_p\to *$, it contains the ordinary localizer generated by that map. In particular all these infinite-loop spaces belong to $\Nul_p$. This supplies the exact consequence of the Kahn--Priddy argument required in REVLOC D.2.2; an infinite-loop section is neither required nor asserted.

For every primary summand $F'$ of every member of $\mathcal C$, realize a sufficiently positive suspension as
$\Sigma^\infty X_{F'}\simeq\Sigma^{r(F')}F'$, where $X_{F'}$ is a finite simply connected pointed complex. Such a realization follows by realizing a finite list of stable attaching maps after a common suspension. If $F'$ is $p$-primary use the countable CW generator $QX_{F'}=\Omega^\infty\Sigma^\infty X_{F'}$; at every other prime use $X_{F'}$ itself. The former is $\Nul_p$-null by the preceding paragraph. For the latter, its $p$-localization is contractible by integral homology and the simply connected Whitehead criterion, and maps into the $p$-local targets factor through that localization. Thus all generators are in $\Nul_p$.

Let $\mathcal A$ be the least weak-equivalence-invariant class of pointed spaces containing these generators and the point, and closed under pointed homotopy colimits, retracts, and all three cofiber implications. Define
\[
                     f\in W\quad\Longleftrightarrow\quad C(Bf)\in\mathcal A.
\]
This is a membership condition for every small-category functor. The proofs in A.5 and A.7 verify every localizer axiom and A.3 verifies all reverse span and cube axioms. In particular no assertion is being made only for open-sieve squares. The countable enlargement class $\mathcal A_c$ of A.7 contains all the chosen generators, which admit countable CW models. It has the same closure properties, and $\mathcal A_c\subseteq\mathcal A$. Minimality therefore gives equality. The generating-set argument in A.7 proves accessibility in ZFC.

Each generator has suspension spectrum in $L$. For $QX_{F'}$, the Snaith splitting expresses that spectrum as a wedge of homotopy orbits of smash powers of $\Sigma^\infty X_{F'}$. All these lie in the tensor ideal $L$. Thus $\mathcal A\subseteq\{X:\Sigma^\infty X\in L\}$ and $T(W)\subseteq L$. Conversely, the degree-one summand of the same splitting exhibits $\Sigma^\infty X_{F'}$ as a retract, and at other primes it is already a generator. Desuspension in the stable category gives every $F'$ and hence $T(W)=L$.

Finally take the chosen nonzero $p$-primary $F$. Its Spanier--Whitehead dual is nonzero, so it has a nonzero stable cohomotopy class in some degree. After a sufficiently large suspension, that class is represented by a map from $\Sigma^rX_F$ to a sphere of odd dimension at least three, localized at $p$. Such a sphere is among the targets for both choices of $p$. Hence $\Sigma^rX_F\notin\Nul_p$ and therefore $X_F\notin\Nul_p$. Since $\mathcal A\subseteq\Nul_p$, we have $X_F\notin\mathcal A$.
Use the named functor $\Delta/\operatorname{Sing}(X_F)\to e$. Its canonical classifying-space comparison with $X_F$ is a weak equivalence; A.2 transfers the asphericity test. Its cofiber spectrum is a suspension of $\Sigma^\infty X_F\in L$, so the functor is in $R(T(W))\setminus W$. For the main counterexample the completely finite named witness is already $P\to e$ of A.5--A.6. The present parameterized realization theorem asserts no unperformed machine enumeration for all finite spectra.

The finite part of $L$ is exactly $\mathcal C$. Indeed a finite spectrum is compact, and the cellular construction of an object of $\loc(\mathcal C)$ uses sums, shifts and successive cofibers of objects of $\mathcal C$. A map from a compact object into that construction factors through a finite subconstruction; applying this to its identity makes it a retract of a finite iterated cofiber of generators. It belongs to $\mathcal C$. The opposite inclusion is immediate. Hopkins--Smith's \cite[Theorem~7]{AHopkinsSmith} classifies each primary finite thick subcategory; the realization proof itself does not require selecting its chromatic index.
\end{proof}
